% Transcription of Cayley's Collected Mathematical Papers, Volume I, pages 101-125
% Paper 12 (conclusion): On the Theory of Determinants
% Paper 13:              On the Theory of Linear Transformations
% Paper 14 (partial):    On Linear Transformations
%
% Method: pdftotext OCR base + scan PNG visual restoration at 200 dpi.
% Source scan: collectedmathema01cayluoft.pdf (Internet Archive). Public domain.

\documentclass[11pt]{article}
\usepackage[a4paper,margin=0.65in]{geometry}
\usepackage[T1]{fontenc}
\usepackage{amsmath,amssymb,amsfonts}
\usepackage{graphicx}
\usepackage{pdfpages}
\usepackage[english,shorthands=off]{babel}
\usepackage{fancyhdr}
\usepackage[bookmarks=false,colorlinks=false]{hyperref}
\usepackage[protrusion=true,expansion=false]{microtype}
\emergencystretch=3em
\tolerance=600
\providecommand{\modd}{\mathrm{mod}}
\providecommand{\Disc}{\mathrm{Disc}}
\providecommand{\canont}{}
\providecommand{\asterism}{*}
\providecommand{\Hessian}{H}
\providecommand{\tome}{}
\providecommand{\page}{p.}
\pagestyle{fancy}\fancyhf{}\fancyhead[L]{Pages 101--125}\fancyhead[R]{Cayley --- Collected Papers, Vol. I}\renewcommand{\headrulewidth}{0.4pt}
\setlength{\headheight}{15pt}

% Cayley's fraktur-style letters used in paper 13 (III, n=4 case)
\newcommand{\fA}{\mathfrak{A}}
\newcommand{\fB}{\mathfrak{B}}
\newcommand{\fC}{\mathfrak{C}}
\newcommand{\fD}{\mathfrak{D}}
\newcommand{\fE}{\mathfrak{E}}
\newcommand{\fF}{\mathfrak{F}}
\newcommand{\fG}{\mathfrak{G}}
\newcommand{\fH}{\mathfrak{H}}
\newcommand{\fI}{\mathfrak{I}}
\newcommand{\fJ}{\mathfrak{J}}
\newcommand{\fK}{\mathfrak{K}}

\begin{document}

% ============================================================
% PAPER 12 (CONCLUSION) -- ON THE THEORY OF DETERMINANTS
% Page 79 (PDF p.101)
% ============================================================

\noindent\textit{[Paper 12 continued from p.~78. Page 79 of the printed volume.]}

\medskip

I proceed to prove a theorem, which may be expressed as follows:
\[
\overset{+}{[A.k.2p]} \cdot \overset{+}{[B.k.2q]} \mid = [\overset{+}{AB} \mid k.2p+2q-2] \tag{27},
\]
where
\[
\overline{AB} \mid_{rs\ldots xy\ldots} = S \cdot A_{rs\ldots l} B_{xy\ldots l} \tag{28},
\]
the number of the symbols $r, s, \ldots$ being obviously $2p-1$, and that of $x, y, \ldots$ being
$2q-1$. The summatory sign $S$ refers to $l$, and denotes the sum of the several terms
corresponding to values of $l$ from $l=1$ to $l=k$. Also the theorem would be equally
true if $l$ had been placed in any position whatever in the series $r, s \ldots l$; and again,
in any position whatever in the series $x, y \ldots l$, instead of at the end of each of these.
With a very slight modification this may be made to suit the case of an odd number
instead of one of the two even numbers $2p,\, 2q$; (in fact, it is only necessary to place
the mark $(+)$ in $\{AB\mid \ldots\}$ over the column corresponding to the marked column in
$\{A\ldots\}$, $\{A\ldots\}$ being the symbol for which the number of columns is odd), but it is
inapplicable where the two numbers are odd. Consider the second side of (27); this
may be expanded in the form
\[
\Sigma \pm_s \ldots \pm_x \pm_y \ldots \overline{AB}\mid_{s_1\ldots x_1y_1\ldots}\,\overline{AB}\mid_{s_2\ldots x_2y_2\ldots}\,\ldots\,\overline{AB}\mid_{s_k\ldots x_ky_k\ldots} \tag{29},
\]
where $\Sigma$ refers to the different quantities $s, \ldots, x, y, \ldots$ as in (11).

Substituting from (28), this becomes
\[
\Sigma \cdot S_{l_1}\ldots S_{l_k}\bigl(\pm_s \pm_x \pm_y \ldots A_{rs_1\ldots l_1}\ldots A_{rs_k\ldots l_k}\,B_{xy_1\ldots l_1}\ldots B_{x_ky_k\ldots l_k}\bigr) \tag{30}.
\]
Effecting the summation with respect to $x, y \ldots$ this becomes
\[
\Sigma \cdot S_{l_1}\ldots S_{l_k} \pm_s \ldots A_{rs_1\ldots l_1}\ldots A_{rs_k\ldots l_k}\left\{
\begin{array}{c}
\overset{+}{B}\,1\,1\,\ldots\,l_1 \\[2pt]
\vdots \\
k\,k\,\ldots\,l_k
\end{array}
\right\} \tag{31},
\]
$S$ now referring to $s, \ldots$ only. The quantity under the sign $S$ vanishes if any two
of the quantities $l$ are equal, and in the contrary case, we have
\[
\left\{
\begin{array}{c}
\overset{+}{B}\,1\,1\,\ldots\,l_1 \\[2pt]
\vdots \\
k\,k\,\ldots\,l_k
\end{array}
\right\}=\pm_l\bigl[\overset{+}{B}.k.2q\bigr] \tag{32},
\]
which reduces the above to
\[
\bigl[\overset{+}{B}.k.2q\bigr]\cdot \Sigma \pm_s \ldots \pm_l A_{rs_1\ldots l_1}\ldots A_{rs_k\ldots l_k} \tag{33},
\]
$\Sigma$ referring to the quantities $s\ldots$, and also to the quantities $l$. And this is evidently
equivalent to
\[
\bigl[\overset{+}{A}.k.2p\bigr]\bigl[\overset{+}{B}.k.2q\bigr] \tag{34},
\]
the theorem to be proved. It is obvious that when $p=1$, $q=1$, the equation (27)
coincides with the theorem $(\Theta)$, quoted in the introduction to this paper.

\bigskip
\hrule
\bigskip

% ============================================================
% PAPER 13 -- ON THE THEORY OF LINEAR TRANSFORMATIONS
% Pages 80--94 (PDF p.102--116)
% ============================================================

\begin{center}
{\Large\bfseries 13.}\\[0.6em]
{\large\bfseries ON THE THEORY OF LINEAR TRANSFORMATIONS.}\\[0.4em]
\textit{[From the Cambridge Mathematical Journal, vol.~iv.~(1845), pp.~193--209.]}
\end{center}

\medskip

The following investigations were suggested to me by a very elegant paper on
the same subject, published in the Journal by Mr Boole. The following remarkable
theorem is there arrived at. If a rational homogeneous function $U$, of the $n^{\text{th}}$ order,
with the $m$ variables $x, y \ldots$, be transformed by linear substitutions into a function
$V$ of the new variables $\xi, \eta \ldots$; if, moreover, $\theta U$ expresses the function of the
coefficients of $U$, which, equated to zero, is the result of the elimination of the
variables from the series of equations $d_x U=0$, $d_y U=0$, \&c., and of course $\theta V$ the
analogous function of the coefficients of $V$: then $\theta V = E^a \cdot \theta U$, where $E$ is the
determinant formed by the coefficients of the equations which connect $x, y \ldots$ with
$\xi, \eta \ldots$, and $a=(n-1)^{m-1}$.\footnote{The value of $a$ was left undetermined, but Mr Boole has since informed me, he was acquainted with
it at the time his paper was written; and has given it in a subsequent paper.} In attempting to demonstrate this very beautiful property,
it occurred to me that it might be generalised by considering for the function $U$, not
a homogeneous function of the $n^{\text{th}}$ order between $m$ variables, but one of the same
order, containing $n$ sets of $m$ variables, and the variables of each set entering
linearly. The form which Mr Boole's theorem thus assumes is $\theta V = E_1^{a_1} \cdot E_2^{a_2} \ldots E_n^{a_n} \cdot \theta U$.
This it was easy to demonstrate would be true, if $\theta U$ satisfied a certain system of
partial differential equations. I imagined at first that these would determine the function
$\theta U$, (supposed, in analogy with Mr Boole's function, to represent the result of the
elimination of the variables from $d_{x_1} U=0$, $d_{y_1} U=0$, \ldots $d_{x_n} U=0$, \&c.): this I afterwards
found was not the case; and thus I was led to a class of functions, including as
a particular case the function $\theta U$, all of them possessed of the same characteristic
property. The system of partial differential equations was without difficulty replaced
by a more fundamental system of equations, upon which, assumed as definitions, the
theory appears to me naturally to depend; and it is this view of it which I intend
partially to develope in the present paper.

\medskip

I have already employed the notation
\[
\left|\,
\begin{array}{cccc}
\alpha, & \beta, & \gamma, & \delta,\,\ldots\\
\alpha', & \beta', & \gamma', & \delta',\,\ldots\\
\alpha'', & \beta'', & \gamma'', & \delta'',\,\ldots\\
\multicolumn{4}{c}{\vdots}
\end{array}
\right| \tag{1}
\]
(where the number of horizontal rows is less than that of vertical ones) to denote
the series of determinants
\[
\left|\,
\begin{array}{ccc}
\alpha, & \beta, & \gamma,\,\ldots\\
\alpha', & \beta', & \gamma',\,\ldots\\
\alpha'', & \beta'', & \gamma'',\,\ldots\\
\multicolumn{3}{c}{\vdots}
\end{array}
\right| \tag{2},
\]
which can be formed out of the above quantities by selecting any system of vertical
rows; these different determinants not being connected together by the sign $+$, or in
any other manner, but being looked upon as perfectly separate.

The fundamental theorem for the multiplication of determinants gives, applied to
these, the formula
\[
\left|\,
\begin{array}{cccc}
A, & B, & C, & D,\,\ldots\\
A', & B', & C', & D',\,\ldots\\
A'', & B'', & C'', & D'',\,\ldots\\
\multicolumn{4}{c}{\vdots}
\end{array}
\right| = E \left|\,
\begin{array}{cccc}
\alpha, & \beta, & \gamma, & \delta,\,\ldots\\
\alpha', & \beta', & \gamma', & \delta',\,\ldots\\
\alpha'', & \beta'', & \gamma'', & \delta'',\,\ldots\\
\multicolumn{4}{c}{\vdots}
\end{array}
\right| \tag{3},
\]
where
\[
\left.
\begin{aligned}
A &= \lambda \alpha + \lambda' \alpha' + \lambda'' \alpha'' + \ldots\\
B &= \lambda \beta + \lambda' \beta' + \lambda'' \beta'' + \ldots\\
A' &= \mu \alpha + \mu' \alpha' + \mu'' \alpha'' + \ldots\\
B' &= \mu \beta + \mu' \beta' + \mu'' \beta'' + \ldots\\
&\ \vdots
\end{aligned}\quad\right\} \tag{4},
\]
\[
E = \left|\,
\begin{array}{cc}
\lambda, & \mu,\,\ldots\\
\lambda', & \mu',\,\ldots\\
\multicolumn{2}{c}{\vdots}
\end{array}
\right| \tag{5},
\]
and the meaning of the equation is, that the terms on the first side are equal, each
to each, to the terms on the second side.

This preliminary theorem being explained, consider a set of arbitrary coefficients,
represented by the general formula
\[
\alpha_{rs\ldots} \tag{6},
\]
in which the number of symbolical letters $r, s, \ldots$ is $n$, and where each of these is
supposed to assume all integer values, from 1 to $m$ inclusively.

Let
\[
\left|\,
\begin{array}{ccc}
1\,s_1' t_1' \ldots, & 1\,s_1'' t_1'' \ldots, & \ldots\\
2\,s_2' t_2' \ldots, & 2\,s_2'' t_2'' \ldots, & \ldots\\
\multicolumn{3}{c}{\vdots}
\end{array}
\right| \tag{7}
\]
represent the whole series, taken in any order, in which the first symbolical letter
is $\alpha$. Similarly,
\[
\left|\,
\begin{array}{ccc}
r_1\,1\,t_1' \ldots, & r_1\,1\,t_1'' \ldots, & \ldots\\
r_2\,2\,t_2' \ldots, & r_2\,2\,t_2'' \ldots, & \ldots\\
\multicolumn{3}{c}{\vdots}
\end{array}
\right| \tag{8},
\]
the whole series of those in which the second symbolical letter is $\alpha$, and so on.

Imagine a function $u$ of the coefficients, which is simultaneously of the forms
\[
\left.
\begin{aligned}
u &= H_p \left|\,
\begin{array}{ccc}
1\,s_1' t_1' \ldots, & 1\,s_1'' t_1'' \ldots, & \ldots\\
2\,s_2' t_2' \ldots, & 2\,s_2'' t_2'' \ldots, & \ldots\\
\multicolumn{3}{c}{\vdots}
\end{array}
\right|\\[8pt]
u &= H_p \left|\,
\begin{array}{ccc}
r_1\,1\,t_1' \ldots, & r_1\,1\,t_1'' \ldots, & \ldots\\
r_2\,2\,t_2' \ldots, & r_2\,2\,t_2'' \ldots, & \ldots\\
\multicolumn{3}{c}{\vdots}
\end{array}
\right|
\end{aligned}\quad\right\} \tag{A};
\]
\&c.; in which $H_p$ denotes a rational homogeneous function of the order $p$. The
function $H$ is not necessarily supposed to be the same in the above equations, and
in point of fact it will not in general be so. The number of equations is of course $=n$.

The function $u$, whose properties we proceed to investigate, may conveniently be
named a ``Hyperdeterminant.'' Any function satisfying any of the equations $(A)$,
without satisfying all of them, will be an ``Incomplete Hyperdeterminant.'' But, considering
in the first place such as are complete---

Let $\overline{rst}\ldots$ be a new set of coefficients connected with the former ones by a system
of equations of the form
\[
\overline{rst}\ldots = \lambda_r{}^1\,1st\ldots + \lambda_r{}^2\,2st\ldots + \ldots + \lambda_r{}^m\,mst\ldots \tag{9},
\]
(where the $r$ in $\lambda_r{}^1\ldots$ is not an exponent, but an affix).

Suppose $\bar{u}$ is the same function of these new coefficients that $u$ was of the former
ones. Then consider the first of the equations $(A)$ and the equation $(3)$, and writing
\[
L = \left|\,
\begin{array}{ccc}
\lambda_1{}^1, & \lambda_2{}^1, & \ldots\\
\lambda_1{}^2, & \lambda_2{}^2, & \ldots\\
\multicolumn{3}{c}{\vdots}
\end{array}
\right| \tag{10},
\]
we have immediately the equation
\[
\bar{u} = L^p \cdot u \tag{11}.
\]

Consider the new set of coefficients
\[
\overline{rst}\ldots = \mu_s{}^1\,r\,1\,t\ldots + \mu_s{}^2\,r\,2\,t\ldots + \ldots + \mu_s{}^m\,r\,m\,t\ldots \tag{12},
\]
and $\bar{u}$ the analogous function of these; then, from the second of the equations $(A)$
and the equation $(3)$, and writing
\[
M = \left|\,
\begin{array}{ccc}
\mu_1{}^1, & \mu_2{}^1, & \ldots\\
\mu_1{}^2, & \mu_2{}^2, & \ldots\\
\multicolumn{3}{c}{\vdots}
\end{array}
\right| \tag{13},
\]
\[
\bar{u} = M^p \cdot u = L^p M^p \cdot u \tag{14}.
\]

In like manner, considering the new coefficients $\overline{rst}\ldots$, where
\[
\overline{rst}\ldots = \nu_t{}^1\,rs\,1\ldots + \nu_t{}^2\,rs\,2\ldots + \ldots + \nu_t{}^m\,rs\,m\ldots \tag{15},
\]
the new function $\bar{u}$, and the quantity $N$ given by
\[
N = \left|\,
\begin{array}{ccc}
\nu_1{}^1, & \nu_2{}^1, & \ldots\\
\nu_1{}^2, & \nu_2{}^2, & \ldots\\
\multicolumn{3}{c}{\vdots}
\end{array}
\right| \tag{16},
\]
we have, as before,
\[
\bar{u} = N^p \cdot u = L^p M^p N^p \cdot u \tag{17},
\]
or
\[
\bar{u} = L^p M^p N^p \cdot u \tag{18};
\]
whence, generally, denoting the last result by $u'$,
\[
u' = L^p M^p N^p \ldots u \tag{B, 1}.
\]

Consider now the function
\[
\Sigma\Sigma\Sigma\ldots\,(rst\ldots \cdot x_r y_s z_t \ldots) \tag{19},
\]
where the $\Sigma$'s refer successively to $r, s, t, \ldots$, and denote summations from 1 to $m$
inclusively. If $u$ be looked upon as a derivative from the above function, we may write
\[
u = \Theta \cdot \Sigma\Sigma\Sigma\ldots\,(rst\ldots \cdot x_r y_s z_t \ldots) \tag{20}.
\]

Assume
\[
\left.
\begin{aligned}
\bar{x}_r &= \lambda_r{}^1\,x_1 + \lambda_r{}^2\,x_2 + \ldots + \lambda_r{}^m\,x_m,\\
\bar{y}_s &= \mu_s{}^1\,y_1 + \mu_s{}^2\,y_2 + \ldots + \mu_s{}^m\,y_m,\\
&\ \vdots
\end{aligned}\quad\right\} \tag{21}.
\]
It is easy to obtain
\[
\Sigma\Sigma\Sigma\ldots\,(\overline{rst}\ldots \cdot x_r y_s z_t \ldots) = \Sigma\Sigma\Sigma\ldots\,(rst\ldots \cdot \bar{x}_r y_s z_t \ldots),\ldots = \Sigma\Sigma\Sigma\ldots\,(rst\ldots \cdot \bar{x}_r \bar{y}_s \bar{z}_t \ldots) \tag{22},
\]
and the formula for $(u)$ becomes
\[
\Theta \cdot \Sigma\Sigma\Sigma\ldots(\overline{rst}\ldots\cdot x_r y_s z_t\ldots) = L^p M^p N^p \ldots \Theta \cdot \Sigma\Sigma\Sigma\ldots(rst\ldots\cdot \bar{x}_r \bar{y}_s \bar{z}_t\ldots) \tag{B, 2}.
\]

Proceeding to obtain the expression of the coefficients $\overline{rst}\ldots$ in terms of the
coefficients $rst\ldots$, we have
\[
\overline{rst}\ldots = \Sigma\Sigma\Sigma\ldots\,(\lambda_r{}^f \mu_s{}^g \nu_t{}^h \ldots fgh\ldots) \tag{C},
\]
where the $\Sigma$'s refer successively to $f, g, h, \ldots$, denoting summations from 1 to $m$
inclusively. Having this equation, it is perhaps as well to retain
\[
u' = L^p M^p N^p \ldots u \tag{B, 1},
\]
instead of (B, 2), that form being principally useful in showing the relation of the
function $u$ to the theory of the transformation of functions.

It may immediately be seen, that in the equations (B), (C) we may, if we please,
omit any number of the marks of variation $(\,)$, omitting at the same time the
corresponding signs $\Sigma$, and the corresponding factors of the series $L, M, N \ldots$

Also, if $u$ be such as only to satisfy some of the equations $(A)$; then, if in the
same formulae we omit the corresponding marks $(\,)$, summatory signs, and terms of the
series $L, M, N \ldots$, the resulting equations are still true.

From the formulae $(A)$ we may obtain the partial differential equations
\[
\Sigma\Sigma\ldots\left(\alpha st\ldots\cdot \frac{d}{d\beta st\ldots}\right)u = 0,\ \text{or}\ pu, \tag{D}
\]
\[
\Sigma\Sigma\ldots\left(r\alpha t\ldots\cdot \frac{d}{dr\beta t\ldots}\right)u = 0,\ \text{or}\ pu,
\]
according as $\alpha$ is not equal, or is equal, to $\beta$;
and so on: the summatory signs referring in every case to those of the series $r, s, t,\ldots$,
which are left variable, and extending from 1 to $m$ inclusively.

To demonstrate this, consider the general form of $u$, as given by the first of the
equations $(A)$. This is evidently composed of a series of terms, each of the form
\[
c PQR \ldots (p\text{ factors}),
\]
in which
\[
P = \left|\,
\begin{array}{ccc}
1\,s_1' t_1' \ldots, & 1\,s_1'' t_1'' \ldots, & \ldots\,(m\text{ terms})\\
\alpha\,s_\alpha' t_\alpha' \ldots, & \alpha\,s_\alpha'' t_\alpha'' \ldots, & \ldots\\
\beta\,s_\beta' t_\beta' \ldots, & \beta\,s_\beta'' t_\beta'' \ldots, & \ldots\\
\multicolumn{3}{c}{\vdots}
\end{array}
\right|
\]
$Q, R, \&\text{c.}$ being of the same form; and we have
\[
\Sigma\Sigma\ldots\left(\alpha st\ldots\cdot\frac{d}{d\beta st\ldots}\right)cPQR\ldots = c\bigl(QR\ldots\,F + PR\ldots Q' + \&\text{c.}\bigr) + \&\text{c.} + \&\text{c.},
\]
and
\[
\Sigma\Sigma\ldots\left(\alpha st\ldots\cdot\frac{d}{d\beta st\ldots}\right)F = \left|\,
\begin{array}{ccc}
1\,s_1' t_1' \ldots, & 1\,s_1'' t_1'' \ldots, & \ldots\,(m\text{ terms})\\
\alpha\,s_\alpha' t_\alpha' \ldots, & \alpha\,s_\alpha'' t_\alpha'' \ldots, & \ldots\\
\alpha\,s_\beta' t_\beta' \ldots, & \alpha\,s_\beta'' t_\beta'' \ldots, & \ldots\\
\multicolumn{3}{c}{\vdots}
\end{array}
\right| = 0;
\]
so that all the terms on the second side of the equation vanish. If, however, $\beta = \alpha$,
whence, on the second side, we have
\[
cQR\ldots P + cPR\ldots Q + \&\text{c.} = p \cdot cPQR\ldots + \&\text{c.} + \&\text{c.} = pU,
\]
or the theorem in question is proved.

In the case of an incomplete hyperdeterminant, the corresponding systems of
equations are of course to be omitted. In every case it is from these equations that
the form of the function $u$ is to be investigated; they entirely replace the system $(A)$.

A very important case of the general theory is, when we suppose the coefficients
$rst\ldots$ to have the property $r's't'\ldots = r''s''t''\ldots$, whenever $r's't'\ldots$ and $r''s''t''\ldots$ denote
the same combination of letters; and also that the coefficients $\lambda$ are equal to the
coefficients $\mu, \nu \ldots$, each to each. In this case the coefficients $\overline{rst}\ldots$ have likewise
the same property, viz. that $\overline{r''s''t''}\ldots = \overline{r's't'}\ldots$, whenever $r's't'\ldots$ and $r''s''t''\ldots$ denote
the same combination of letters.

The equations $(B, 1)$, $(B, 2)$ become in this case
\[
u' = L^{\rho p}\cdot u \tag{B, 3},
\]
\[
\Theta\cdot\Sigma\Sigma\Sigma\ldots\left[\frac{[\alpha]^\rho}{[\alpha']^\rho [\beta']^\rho \ldots}\overline{rst}\ldots \bar{x}_r \bar{y}_s \bar{z}_t \ldots\right] = L^{\rho p}\,\Theta\cdot\Sigma\Sigma\Sigma\ldots\left[\frac{[\alpha]^\rho}{[\alpha']^\rho [\beta']^\rho \ldots}\,rst\ldots \cdot x_r y_s z_t \ldots\right] \tag{B, 4},
\]
where only different combinations of values are to be taken for $r, s, t, \ldots$ and
$\alpha, \beta, \ldots$ express how often the same number occurs in the series. In the equation
$(C)$, $\mu, \nu$ must be replaced by $\lambda$, the equations $(D)$ are no longer satisfied; the
equations $(A)$ reduce themselves to a single one, (so that there can be no question
here of incomplete hyperdeterminants): but this is no longer sufficient to determine
the function sought after. For this reason, the particular case, treated separately,
would be far more difficult than the general one; but the formulae of the general case
being first established, these apply immediately to the particular one\footnote{See concluding paragraph of this paper.}. The case in
question may be defined as that of symmetrical hyperdeterminants, (a denomination
already adopted for ordinary determinants). It would be easily seen what on the same
principle would be meant by partially symmetrical hyperdeterminants.

I have not yet succeeded in obtaining the general expression of a hyperdeterminant;
the only cases in which I can do so are the following: I. $p=1$, $n$ even, (if $n$ be odd,
there only exist incomplete hyperdeterminants). II. $p=2$, $m=2$, $n$ even. III. $p=3$,
$m=2$, $n=4$.

I. The first case is, in fact, that of the functions considered at the termination of
a paper in the Cambridge Philosophical Transactions, vol.~viii.~part i.~[12]; though at
that time I was quite unacquainted with the general theory.

Using the notation there employed, we have
\[
u = \left\{
\begin{array}{c}
\overset{+}{1\,1\,\ldots\,(n)}\\
2\,2\\
\vdots\\
mm
\end{array}
\right\},
\]
a complete hyperdeterminant when $n$ is even; and when $n$ is odd the functions
\[
\left\{
\begin{array}{c}
\overset{+}{1\,1\,\ldots\,(n)}\\
2\,2\\
\vdots\\
mm
\end{array}
\right\},\qquad
\left\{
\begin{array}{c}
\overset{+}{1\,1\,\ldots\,(n)}\\
2\,2\\
\vdots\\
mm
\end{array}
\right\}
\]
are each of them incomplete hyperdeterminants.

(A) In the case of $n=2$, the complete hyperdeterminant is simply the ordinary
determinant
\[
\left|\,
\begin{array}{cccc}
1\,1, & 1\,2, & \ldots & 1m\\
2\,1, & 2\,2, & \ldots & 2m\\
\multicolumn{4}{c}{\vdots}\\
m1, & m2, & \ldots & mm
\end{array}
\right|.
\]
Stating the general conclusion as applied to this case, which is a very well known one,

``If the function
\[
U = \Sigma\Sigma\,(rs \cdot x_r y_s)
\]
be transformed into a similar function
\[
\Sigma\Sigma\,(\overline{rs}\cdot \bar{x}_r \bar{y}_s).
\]
by means of the substitutions
\[
\bar{x}_r = \lambda_r{}^1\,x_1 + \lambda_r{}^2\,x_2 + \ldots + \lambda_r{}^m\,x_m,
\]
\[
\bar{y}_s = \mu_s{}^1\,y_1 + \mu_s{}^2\,y_2 + \ldots + \mu_s{}^m\,y_m;
\]
then
\[
\left|\,
\begin{array}{cccc}
\overline{1\,1}, & \overline{1\,2}, & \ldots\\
\overline{2\,1}, & \overline{2\,2}, & \\
\multicolumn{3}{c}{\vdots}
\end{array}
\right| =
\left|\,
\begin{array}{cc}
\lambda_1{}', & \lambda_2{}',\,\ldots\\
\lambda_1{}'', & \lambda_2{}'',\,\ldots\\
\multicolumn{2}{c}{\vdots}
\end{array}
\right|
\left|\,
\begin{array}{cc}
\mu_1{}', & \mu_2{}',\,\ldots\\
\mu_1{}'', & \mu_2{}'',\,\ldots\\
\multicolumn{2}{c}{\vdots}
\end{array}
\right|
\left|\,
\begin{array}{ccc}
1\,1, & 1\,2, & \ldots\\
2\,1, & 2\,2, & \\
\multicolumn{3}{c}{\vdots}
\end{array}
\right|.
\]''

Also, by what has preceded,
\[
H = \Sigma\Sigma\,(\Theta\cdot fg \cdot \lambda_r{}^f \mu_s{}^g);
\]
so that the theorem is easily seen to amount to the following one---``If the terms of
a determinant of the $m^{\text{th}}$ order be of the form $\Sigma_r \Sigma_s\,(rs \cdot x_r \mu_s)$, $r, s$ extending as
before, from 1 to $m$ inclusively, the determinant itself is the product of three determinants;
the first formed with the coefficients $rs$, the second with the quantities $x$,
and the third with the quantities $y$.''

In a following number of the Journal I shall prove, and apply to the theories of
Maxima and Minima and of Spherical Coordinates, (I may just mention having obtained,
in an elegant form, the formulae for transforming from one oblique set of coordinates
to another oblique one) the more general theorem,

``If $k$ be the order of the determinant formed as above, the determinant itself
is a quadratic function, its coefficients being determinants formed with the coefficients
$rs$, its variables being determinants formed respectively with the variables $x$ and the
variables $y$; and the number of variables in each set being the number of combinations
of $k$ things out of $m$, $(=1$ if $k=m$; if $k > m$ the determinant vanishes).''

I shall give in the same paper the demonstration of a very beautiful theorem,
rather relating, however, to determinants than to quadratic functions, proved by Hesse
in a Memoir in Crelle's Journal, vol.~xx., ``De curvis et superficiebus secundi ordinis;''
and from which he has deduced the most interesting geometrical results. Another
Memoir, by the same author, Crelle, vol.~XXVIII., ``Ueber die Elimination der Variabeln
aus drei algebraischen Gleichungen vom zweiter Grade mit zwei Variabeln,'' though
relating in point of fact rather to functions of the third order, contains some most
important results. A few theorems on quadratic functions, belonging, however, to a
different part of the subject, will be found in my paper already quoted in the Cambridge
Philosophical Transactions [12]; and likewise in a paper in the Journal, Chapters in
the Algebraical Geometry of $n$ dimensions [11].

I shall, just before concluding this case, write down the particular formula corresponding
to three variables, and for the symmetrical case. It is, as is well known, the
theorem,

``If
\[
U = Ax^2 + By^2 + Cz^2 + 2Fyz + 2Gxz + 2Hxy
\]
be transformed into
\[
\fA \xi^2 + \fB \eta^2 + \fC \zeta^2 + 2\fF \eta\zeta + 2\fG \xi\zeta + 2\fH \xi\eta
\]
by means of
\[
\begin{aligned}
x &= \alpha\xi + \beta\eta + \gamma\zeta,\\
y &= \alpha'\xi + \beta'\eta + \gamma'\zeta,\\
z &= \alpha''\xi + \beta''\eta + \gamma''\zeta,
\end{aligned}
\]
then
\[
(\fA\fB\fC - \fA\fF^2 - \fB\fG^2 - \fC\fH^2 + 2\fF\fG\fH) =
\]
\[
(\alpha\beta'\gamma'' - \alpha\beta''\gamma' + \alpha'\beta''\gamma - \alpha'\beta\gamma'' + \alpha''\beta\gamma' - \alpha''\beta'\gamma)^2\,(ABC - AF^2 - BG^2 - CH^2 + 2FGH).''
\]

(B) Let $n=3$, and for greater simplicity $m=2$; write
\[
\begin{aligned}
a &= 111, & e &= 112,\\
b &= 211, & f &= 212,\\
c &= 121, & g &= 122,\\
d &= 221, & h &= 222,
\end{aligned}
\]
so that
\[
U = a\,x_1 y_1 z_1 + b\,x_2 y_1 z_1 + c\,x_1 y_2 z_1 + d\,x_2 y_2 z_1 + e\,x_1 y_1 z_2 + f\,x_2 y_1 z_2 + g\,x_1 y_2 z_2 + h\,x_2 y_2 z_2.
\]
There is no complete hyperdeterminant (i.e. for $p=1$), and the incomplete ones are
\[
\begin{aligned}
ah - bg - cf + de &= u_1,\ \text{suppose,}\\
ah - de - bg + cf &= u_2,\\
ah - cf - de + bg &= u_3.
\end{aligned}
\]

Thus, suppose the transforming equations are
\[
\begin{aligned}
\bar{x}_1 &= \lambda_1{}'\,x_1 + \lambda_1{}''\,x_2,\\
\bar{x}_2 &= \lambda_2{}'\,x_1 + \lambda_2{}''\,x_2;\\
\bar{y}_1 &= \mu_1{}'\,y_1 + \mu_1{}''\,y_2,\\
\bar{y}_2 &= \mu_2{}'\,y_1 + \mu_2{}''\,y_2;\\
\bar{z}_1 &= \nu_1{}'\,z_1 + \nu_1{}''\,z_2,\\
\bar{z}_2 &= \nu_2{}'\,z_1 + \nu_2{}''\,z_2;
\end{aligned}
\]
then
\[
\begin{aligned}
\bar{u}_1 &= (\mu_1'\mu_2'' - \mu_2'\mu_1'')\,(\nu_1'\nu_2'' - \nu_2'\nu_1'')\,u_1,\ \text{where } y, z \text{ are changed,}\\
\bar{u}_2 &= (\nu_1'\nu_2'' - \nu_2'\nu_1'')\,(\lambda_1'\lambda_2'' - \lambda_2'\lambda_1'')\,u_2,\quad ''\ z, x\ ''\\
\bar{u}_3 &= (\lambda_1'\lambda_2'' - \lambda_2'\lambda_1'')\,(\mu_1'\mu_2'' - \mu_2'\mu_1'')\,u_3,\quad ''\ x, y\ ''
\end{aligned}
\]
We might also have assumed
\[
\begin{aligned}
u_1 &= ad - bc,\ \text{or}\ eh - gf,\\
u_2 &= af - be,\ \text{or}\ ch - dg,\\
u_3 &= ag - ce,\ \text{or}\ bh - df,
\end{aligned}
\]
but these are ordinary determinants.

(C) $n=4$, $m=2$.
\[
\begin{aligned}
a &= 1111, & i &= 1112,\\
b &= 2111, & j &= 2112,\\
c &= 1211, & k &= 1212,\\
d &= 2211, & l &= 2212,\\
e &= 1121, & m &= 1122,\\
f &= 2121, & n &= 2122,\\
g &= 1221, & o &= 1222,\\
h &= 2221, & p &= 2222.
\end{aligned}
\]
\[
\begin{aligned}
U &= a\,x_1 y_1 z_1 w_1 + b\,x_2 y_1 z_1 w_1 + c\,x_1 y_2 z_1 w_1 + d\,x_2 y_2 z_1 w_1\\
&\quad + e\,x_1 y_1 z_2 w_1 + f\,x_2 y_1 z_2 w_1 + g\,x_1 y_2 z_2 w_1 + h\,x_2 y_2 z_2 w_1\\
&\quad + i\,x_1 y_1 z_1 w_2 + j\,x_2 y_1 z_1 w_2 + k\,x_1 y_2 z_1 w_2 + l\,x_2 y_2 z_1 w_2\\
&\quad + m\,x_1 y_1 z_2 w_2 + n\,x_2 y_1 z_2 w_2 + o\,x_1 y_2 z_2 w_2 + p\,x_2 y_2 z_2 w_2,
\end{aligned}
\]
we have
\[
u = ap - bo - cn + dm - el + fk + gj - hi;
\]
so that, with the same sets of transforming equations as above, and the additional one,
\[
\begin{aligned}
\bar{w}_1 &= \rho_1'\,w_1 + \rho_1''\,w_2,\\
\bar{w}_2 &= \rho_2'\,w_1 + \rho_2''\,w_2,
\end{aligned}
\]
we have
\[
\bar{u} = (\lambda_1'\lambda_2'' - \lambda_2'\lambda_1'')(\mu_1'\mu_2'' - \mu_2'\mu_1'')(\nu_1'\nu_2'' - \nu_2'\nu_1'')(\rho_1'\rho_2'' - \rho_2'\rho_1'')\cdot u;
\]
this is important when viewed in reference to a result which will presently be obtained.

If we take the symmetrical case, we have
\[
U = \alpha x^4 + 4\beta x^3 y + 6\gamma x^2 y^2 + 4\delta x y^3 + \epsilon y^4,
\]
which is transformed into
\[
U' = \alpha' x'^4 + 4\beta' x'^3 y' + 6\gamma' x'^2 y'^2 + 4\delta' x' y'^3 + \epsilon' y'^4,
\]
by means of
\[
\begin{aligned}
x &= \lambda x' + \mu y',\\
y &= \lambda' x' + \mu' y';
\end{aligned}
\]
if
\[
u = \alpha\epsilon - 4\beta\delta + 3\gamma^2,
\]
\[
u' = \alpha'\epsilon' - 4\beta'\delta' + 3\gamma'^2,
\]
we have
\[
u' = (\lambda\mu' - \lambda'\mu)^4 \cdot u.
\]

\bigskip
II. Where $p=2$, $m=2$, $n$ is odd.

The expression
\[
u = \left|\,
\begin{array}{c}
\overset{+}{(1111\ldots(n))}\\
\hline
1222\\
\hline
1111\ldots(n)\\
\hline
2222
\end{array}
\right|^2 + \left|\,
\begin{array}{c}
\overset{+}{(1111\ldots(n))}\\
\hline
2222\\
\hline
2111\ldots\\
\hline
2222
\end{array}
\right|
\]
is a complete hyperdeterminant; and that over whichever row the mark $(\dagger)$ of non-permutation
is placed. The different expressions so obtained are not, however, all of them
independent functions: thus, in the following example, where $n=3$, the three functions
are absolutely identical.

(A). $n=3$, notation as in I.~(B).
\[
u = a^2 h^2 + b^2 g^2 + c^2 f^2 + d^2 e^2 - 2ahbg - 2ahcf - 2ahde - 2bgcf - 2bgde - 2cfde + 4adfg + 4bech,
\]
and then
\[
\bar{u} = (\lambda_1'\lambda_2'' - \lambda_2'\lambda_1'')^2(\mu_1'\mu_2'' - \mu_2'\mu_1'')^2(\nu_1'\nu_2'' - \nu_2'\nu_1'')^2\,u.
\]
This is in many respects an interesting example. We see that the function $u$
may be expressed in the three following forms:
\[
u = (ah - bg - cf + de)^2 + 4\,(ad - bc)(fg - eh) \tag{1},
\]
\[
u = (ah - bg - de + cf)^2 + 4\,(af - be)(dg - ch) \tag{2},
\]
\[
u = (ah - cf - de + bg)^2 + 4\,(ag - ce)(df - bh) \tag{3},
\]
which are indeed the direct results of the general form above given, the sign $(\dagger)$
being placed in succession over the different columns: and the three forms, as just
remarked, are in this case identical.

We see from the first of these that $u$ is of the second or third, from the second
that $u$ is of the first or third, from the third that $u$ is of the first or second of the
three following forms:
\[
u = H_2\!\left|\,\begin{array}{cc} a, b, c, d \\ e, f, g, h \end{array}\right|,\quad
u = H_2\!\left|\,\begin{array}{cc} a, b, e, f \\ c, d, g, h \end{array}\right|,\quad
u = H_2\!\left|\,\begin{array}{cc} a, c, e, g \\ b, d, f, h \end{array}\right|,
\]
which is as it should be.

The following is a singular property of $u$:

Let
\[
a' = \tfrac{1}{2}\,\frac{du}{da},\quad b' = \tfrac{1}{2}\,\frac{du}{db},\quad \ldots,\quad h' = \tfrac{1}{2}\,\frac{du}{dh};
\]
then, $u'$ being the same function of these new coefficients that $u$ is of the former ones,
\[
u' = u^3.
\]

To prove this, write
\[
p = ah - bg - cf + de,\quad q = (ad - bc),\quad r = eh - fg;
\]
\[
\begin{aligned}
a_1 &= ap - 2qe, & e_1 &= -2ra + pe,\\
b_1 &= bp - 2qf, & f_1 &= -2rb + pf,\\
c_1 &= cp - 2qg, & g_1 &= -2rc + pg,\\
d_1 &= dp - 2qh, & h_1 &= -2rd + ph;
\end{aligned}
\]
we have, as a particular case of the general formula just obtained,
\[
u_1 = (p^2 - 4qr)^2\,u = u_1{}^2\cdot u,\ =u^3.
\]
Also
\[
\begin{aligned}
a_1 &= h', & e_1 &= d',\\
b_1 &= -g', & f_1 &= -c',\\
c_1 &= -f', & g_1 &= -b',\\
d_1 &= e', & h_1 &= a';
\end{aligned}
\]
whence $u_1 = u'$, that is $u' = u^3$.

There is no difficulty in showing also, that if $a'', b'', \ldots h''$ are derived from
$a', b', \ldots h'$, as these are from $a, b, \ldots h$, then
\[
a'' = u^2\,a,\quad b'' = u^2\,b,\,\ldots\,h'' = u^2\,h.
\]
The particular case of this theorem, which corresponds to symmetrical values of the
coefficients, is given by M. Eisenstein, Crelle, vol.~xxvii.~[1844], as a corollary to his
researches on the cubic forms of numbers.

Considering this symmetrical case
\[
U = \alpha x^3 + 3\beta x^2 y + 3\gamma x y^2 + \delta y^3,
\]
\[
u = \alpha^2\delta^2 - 6\alpha\delta\beta\gamma - 3\beta^2\gamma^2 + 4\beta^3\delta + 4\alpha\gamma^3,
\]
so that if $U$ be transformed into
\[
U' = \alpha' x'^3 + 3\beta' x'^2 y' + 3\gamma' x' y'^2 + \delta' y'^3,
\]
by means of
\[
\begin{aligned}
x &= \lambda x' + \mu y',\\
y &= \lambda' x' + \mu' y',
\end{aligned}
\]
then if
\[
u' = \alpha'^2\delta'^2 - 6\alpha'\delta'\beta'\gamma' - 3\beta'^2\gamma'^2 + 4\beta'^3\delta' + 4\alpha'\gamma'^3,
\]
we have
\[
u' = (\lambda\mu' - \lambda'\mu)^6 \cdot u.
\]

\bigskip
III.\quad $p=3$,\ $m=2$,\ $n=4$.

Notation as in I.~(C),
\[
u = A\,(\fA + 3\fB + 3\fC + 6\fD + 6\fE) - B\,(\fF + \fG - 3\fH - 3\fI + 2\fJ + 3\fK),
\]
where $A, B$ are arbitrary constants, and $\fA, \fB,$ \&c.\,\ldots $\fK$ are functions of the coefficients,
given as follows:---
\[
\begin{aligned}
\fA &= -a^2 p^2 io + b^2 o^2 ap + c^2 n^2 dm - d^2 m^2 hi + e^2 l^2 fk - f^2 k^2 el - g^2 j^2 hi + h^2 i^2 gj\\
&\quad - a^2 p^2 cn + b^2 o^2 dm + c^2 n^2 ap - d^2 m^2 bo + e^2 l^2 gj - f^2 k^2 hi - g^2 j^2 el + h^2 i^2 fk\\
&\quad - a^2 p^2 el + b^2 o^2 fk + c^2 n^2 gj - d^2 m^2 hi + e^2 l^2 ap - f^2 k^2 bo - g^2 j^2 cn + h^2 i^2 dm\\
&\quad - a^2 p^2 hi + b^2 o^2 gj + c^2 n^2 fk - d^2 m^2 el + e^2 l^2 dm - f^2 k^2 cn - g^2 j^2 bo + h^2 i^2 ap.
\end{aligned}
\]
\[
\begin{aligned}
\fB &= +a^2 p^2 dm - b^2 o^2 cn - c^2 n^2 bo + d^2 m^2 ap - e^2 l^2 hi + f^2 k^2 gj + g^2 j^2 fk - h^2 i^2 el\\
&\quad + a^2 p^2 fk - b^2 o^2 el - c^2 n^2 hi + d^2 m^2 gj - e^2 l^2 bo + f^2 k^2 ap + g^2 j^2 dm - h^2 i^2 cn\\
&\quad + a^2 p^2 gj - b^2 o^2 hi - c^2 n^2 el + d^2 m^2 fk - e^2 l^2 cn + f^2 k^2 dm + g^2 j^2 ap - h^2 i^2 bo.
\end{aligned}
\]
\[
\begin{aligned}
\fC &= ap\,bocn - ap\,bodm - ap\,cndm + bocndm - elfkgj + elfkhi + elgjhi - higjfk\\
&\quad + ap\,boel - ap\,bojk - cndmgj + dmcnhi - elfkap + elfkbo + gjhicn - higjdm\\
&\quad - ap\,bogj + ap\,bohi + cndmel - dmcnkf + elfkcn - elfkdm - gjhiap + higjbo\\
&\quad + ap\,cnel - bodmfk - cnapgj + dmboki - elgjap + fkhibo + gjelcn - hifkdm\\
&\quad - ap\,cnfk + bodmel + cnaphi - dmboel + elgjbo - fkhiap - gjeldm + hifkcn\\
&\quad - ap\,dmel + bocnkf + cmbogj - dmaphi + elhiap - fkgjbo - gjfkcn + hidmel.
\end{aligned}
\]
\[
\fD = ap\,dmjk - bocnel - bocnik + dmapgj - elhibo + fkapgi + gjfkdm - hielcn.
\]
\[
\begin{aligned}
\fE &= a^2 phjo - b^2 goip - c^2 hifn + d^2 emkn - e^2 dlkn + f^2 ckml + g^2 bjpi - h^2 aijo\\
&\quad - i^2 phbg + j^2 ogah + k^2 nfde - l^2 mecf + m^2 ldcf - n^2 ckde - o^2 bjah + p^2 aibg\\
&\quad + a^2 phkn - b^2 golm - c^2 nfip + d^2 emjo - e^2 dljo + f^2 ckip + g^2 bjpl - h^2 imk\\
&\quad - i^2 phcf + j^2 ogde + k^2 nfah - l^2 mebg + m^2 ldbg - n^2 ckah - o^2 bjde + p^2 aicf\\
&\quad + a^2 phlm - b^2 gokn - c^2 hifjo + d^2 emip - e^2 dlpi + f^2 ckjo + g^2 bjnk - h^2 aiml\\
&\quad - i^2 phde + j^2 ogcf + k^2 nfbg - l^2 meah + m^2 ldah - n^2 ckbg - o^2 bjcf + p^2 aide\\
&\quad + a^2 pdno - b^2 cmpo - c^2 bpmn + d^2 aomn - e^2 hjkl + f^2 gilk + g^2 flij - h^2 ekji\\
&\quad - i^2 lfgh + j^2 kehg + k^2 jhef - l^2 igfe + m^2 pbcd - n^2 aodc - o^2 ndab + p^2 mbca\\
&\quad + a^2 plng - b^2 okmo - c^2 ejpn + d^2 fiom - e^2 cpjl + f^2 doik + g^2 anlj - h^2 bmki\\
&\quad - i^2 odfh + j^2 pceg + k^2 mhbf - l^2 nage + m^2 khbd - n^2 lgac - o^2 ifbd + p^2 heca\\
&\quad + a^2 plfo - b^2 ekpo - c^2 hjmn + d^2 gimn - e^2 bpkl + f^2 aolk + g^2 dnij - h^2 cmji\\
&\quad - i^2 hdgh + j^2 nchg + k^2 pbef - l^2 oafe + m^2 jhcd - n^2 lgcd - o^2 ifab + p^2 keba.
\end{aligned}
\]
{\small
\[
\begin{aligned}
\fF &= ap\,bgkn - boalhm - cndiep + dmcjfo\\
&\quad - elfocj + fkpied + gjhinla - higjnhk\\
&\quad - ihjocf + gjidep + kflamh - lekbng\\
&\quad + mdnhkg - ncmhal - obpied + paofjc\\
&\quad + ap\,bglm - boahkn - cndejo + dmcfip\\
&\quad - elcfip + fkedoj + gjhakn - hibgml\\
&\quad - ihjode + gijpcf + kflmbg - leknah\\
&\quad + mdknah - ncmlbg - obpicf + paojde\\
&\quad + ap\,cflg - bojehk - onjehk + dmilgf\\
&\quad - elmpbc + fkadno + gjadno - hipmcb\\
&\quad - ihadno + gjbmcp + kfcbpm - eladno\\
&\quad + mdjehk - ncilfg - obilfg + pahejk\\
&\quad + ap\,cfhn - bodhie - cnahjo + dmbgip\\
&\quad - elbgip + fkahjo + gjednk - hicfml\\
&\quad - ihknde + mdahjo + kfipbg - bocfml\\
&\quad + gjcfml - ncpigb - leahjo + paknde\\
&\quad + ap\,idng - bojcmh - cnbkpe + dmaflo\\
&\quad - elmhjc - fkidng + gjaflo - ihbkpe\\
&\quad - ihaflo + gjbkpe + fkcjhm - elidng\\
&\quad + mdbkpe - cnaflo - boidng + pahmcj\\
&\quad + ap\,idfo - bojcep - cnbkhm + dmialng\\
&\quad - elmnbk - fkalng + gjidfo - ihpecj\\
&\quad - ihalng + gjkbmh + fkcjpe - elidfo\\
&\quad + mdjcep - ncidfo - obalvg + pahmbk.
\end{aligned}
\]
}
\[
\begin{aligned}
\fG &= a^2 hord - b^2 pgmk - c^2 pfnj + d^2 onie - e^2 dpkj + f^2 lkco + g^2 lni - h^2 amkj\\
&\quad - i^2 pdfg + j^2 oech + k^2 nbeh - l^2 mafg + m^2 lch - n^2 kadg - o^2 jadf + p^2 ibec :
\end{aligned}
\]
we have, as usual,
\[
\bar{u} = (\lambda_1'\lambda_2'' - \lambda_2'\lambda_1'')^3(\mu_1'\mu_2'' - \mu_2'\mu_1'')^3(\nu_1'\nu_2'' - \nu_2'\nu_1'')^3(\rho_1'\rho_2'' - \rho_2'\rho_1'')^3\cdot u.
\]
Particular forms of $U$ are
\smallskip

\noindent $A=1$, $B=0$:
\[
u = \fA + 3\fB + 3\fC + 6\fD + 6\fE = (ap - bo - cn + dm - el + fk + gj - hi)^2,\ =v\ \text{suppose}.
\]

\noindent $A=1$, $B=9$:
\[
u = \fA + 3\fB - 6\fC - 3\fD + 3\fE\fF + 9\fG - 18\fH - 27\fK,\ =\theta U\ \text{suppose,}
\]
where $\theta U=0$ is the result of the elimination of the variables from the equations
$d_{x_1} U=0$, $d_{y_1} U=0$, $d_{z_1} U=0$, $d_{w_1} U=0$, $d_{x_2} U=0$, $d_{y_2} U=0$, $d_{z_2} U=0$, $d_{w_2} U=0$. In fact, by
an investigation similar to Mr Boole's, applied to a function such as $U$, it is shown
that $\theta U$ has the characteristic property of the function $u$: also in the present case
$u$ is the most general function of its kind, so that $\theta U$ is obtained from $U$ by
properly determining the constants. This has been effected by comparing the value of
$u$, in the symmetrical case, with the value of $\theta U$, in the same case, the expanded
expression of which is given by Mr Boole in the Journal, vol.~iv.~p.~169. Assuming
$A=1$, the result was $B=9$. [Incorrect; the result of the elimination is not $\theta U = 0$,
but an equation of a higher degree.]

The general form of $u$ now becomes
\[
u = \alpha v^2 + \beta\,\theta U
\]
in which $\alpha, \beta,$ are indeterminate.

We have
\[
\bar{u} = \alpha \bar{v}^2 + \beta\,\theta\bar{U} = M\,(\alpha v^2 + \beta\,\theta U),
\]
where
\[
M = (\lambda_1'\lambda_2'' - \lambda_2'\lambda_1'')^4(\mu_1'\mu_2'' - \mu_2'\mu_1'')^4(\nu_1'\nu_2'' - \nu_2'\nu_1'')^4(\rho_1'\rho_2'' - \rho_2'\rho_1'')^4,
\]
and thence
\[
\bar{v}^2 = M\,v^2,
\]
which coincides with a previous formula, and
\[
\theta\bar{U} = M\,\theta U:
\]
whence, eliminating $M$,
\[
\frac{\theta\bar{U}}{\bar{v}^2} = \frac{\theta U}{v^2}
\]
an equation which is remarkable as containing only the constants of $U$ and $\bar{U}$: it is
an equation of condition which must exist among the constants of $\bar{U}$ in order that
this function may be derivable by linear substitutions from $U$.

In the symmetrical case, or where
\[
U = \alpha x^4 + 4\beta x^3 y + 6\gamma x^2 y^2 + 4\delta x y^3 + \epsilon y^4,
\]
it has been already seen that $v$ is given by
\[
v = \alpha\epsilon - 4\beta\delta + 3\gamma^2.
\]

Proceeding to form $\theta U$, we have
\[
\begin{aligned}
\fA &= \alpha^2\epsilon^2 - 4\beta^3\delta^3 + 3\gamma^4,\\
\fB &= 4\,(\alpha\epsilon^2\delta^2 - \alpha^2\epsilon^2\delta + 3\gamma^2\beta^2\delta^2 - 4\beta^3\delta^3),\\
\fC &= 3\,(\alpha^2 \epsilon^2 \gamma^2 + 2\gamma^4 + \alpha\epsilon\gamma^4 - 4\beta^2 \delta^2),\\
\fD &= 6\,(\alpha\epsilon^2 \delta^2 - 2\alpha\epsilon\beta\delta\gamma^2 + 3\beta^2\gamma^2\delta^2 - 2\beta\delta\gamma^4),\\
\fE &= 3\alpha\epsilon\gamma^4 - 4\beta^3\delta^3 + \gamma^4,\\
\fF &= 6\,(\alpha^2\gamma\epsilon + \epsilon^2\beta^2\gamma - 2\beta^3\epsilon\gamma\delta - 2\delta^3\alpha\beta\gamma - 4\beta^2\gamma^2\delta^2 + 4\,\beta\delta\gamma^2 + \gamma^3\beta^2\epsilon + \gamma\alpha\delta^2),\\
\fG &= 12\,(\alpha\epsilon\beta\delta\gamma^2 - \alpha\gamma^3\delta^2 - \epsilon\gamma^3\beta^2 + \gamma^2\alpha\delta^2 + \gamma\epsilon\beta^3 + \beta\delta\gamma^4 - 2\beta^2\gamma^2\delta^2),\\
\fK &= (\alpha^2\delta^4 + \epsilon^2\beta^4 - 4\beta^3\epsilon\gamma^3 - 4\alpha\gamma\beta^3 + 6\beta^2\gamma^2\delta^2),
\end{aligned}
\]
and these values give
\[
\begin{aligned}
\theta U &= \alpha^2\epsilon^4 - 6\alpha\beta^2\delta^2\epsilon - 12\alpha^2\beta\delta\epsilon^2 - 18\alpha^2\gamma^2\epsilon^2 - 27\alpha^2\delta^4 - 27\beta^4\epsilon^2 + 36\beta^2\gamma\delta^2\epsilon\\
&\quad + 54\alpha\beta\gamma\delta^2\epsilon + 54\alpha\beta^2\gamma\epsilon^2 - 54\alpha\gamma^3\delta^2 - 54\beta^3\gamma^2\epsilon - 64\beta^3\delta^3 + 81\alpha^2\delta\epsilon\\
&\quad + 108\alpha\gamma\beta\delta^3 + 108\beta^2\gamma\delta\epsilon - 180\alpha\beta\gamma^2\delta\epsilon,
\end{aligned}
\]
so that this function, divided by $(\alpha\epsilon - 4\beta\delta + 3\gamma^2)^3$, is invariable for all functions of the
fourth order which can be deduced one from the other by linear substitutions. The
function $\alpha\epsilon - 4\beta\delta + 3\gamma^2$ occurs in other investigations: I have met with it in a problem
relating to a homogeneous function of two variables, of any order whatever, $\alpha, \beta, \gamma, \delta, \epsilon$
signifying the fourth differential coefficients of the function. But this is only remotely
connected with the present subject.

Since writing the above, Mr Boole has pointed out to me that in the transformation
of a function of the fourth order of the form $a x^4 + 4 b x^3 y + 6 c x^2 y^2 + 4 d x y^3 + e y^4$,---
besides his function $\theta u$, and my quadratic function $ae - 4 b d + 3 c^2$,---there exists a function
of the third order $ace - b^2 e - a d^2 - c^3 + 2 b d c$, possessing precisely the same characteristic
property, and that, moreover, the function $\theta u$ may be reduced to the form
\[
(ae - 4bd + 3c^2)^3 - 27\,(ace - ad^2 - eb^2 - c^3 + 2bdc)^2;
\]
the latter part of which was verified by trial; the former he has demonstrated in a
manner which, though very elegant, does not appear to be the most direct which the
theorem admits of. In fact, it may be obtained by a method just hinted at by
Mr Boole, in his earliest paper on the subject, Mathematical Journal, vol.~II.~p.~70.
The equations $d_x^2 u = 0$, $d_x d_y u = 0$, $d_y^2 u = 0$, imply the corresponding equations for the
transformed function: from these equations we might obtain two relations between the
coefficients, which, in the case of a function of the fourth order, are of the orders 3
and 4 respectively: these imply the corresponding relations between the coefficients of
the transformed function. Let $A=0$, $B=0$, $A'=0$, $B'=0$, represent these equations;
then, since $A=0$, $B=0$, imply $A'=0$, we must have $A' = \Lambda A + M B$, $\Lambda, M,$ being
functions of $\lambda, \lambda', \mu, $\&c.~$\mu'$: but $B$ being of the fourth order, while $A, A'$ are only of
the third order in the coefficients of $u$, it is evident that the term $MB$ must disappear,
or that the equation is of the form $A' = \Lambda A$. The function $A$ is obviously the function
which, equated to zero, would be the result of the elimination of $x^2, xy, y^2$, considered
as independent quantities from the equations $a x^2 + 2 b x y + c y^2 = 0$, $b x^2 + 2 c x y + d y^2 = 0$,
$c x^2 + 2 d x y + e y^2 = 0$, viz. the function given above. Hence the two functions on which
the linear transformation of functions of the fourth order ultimately depend are the
very simple ones
\[
ae - 4bd + 3c^2,\quad ace - ad^2 - eb^2 - c^3 + 2bdc,
\]
the function of the sixth order being merely a derivative from these. The above
method may easily be extended: thus for instance, in the transformation of functions of
any even order, I am in possession of several of the transforming functions; that of the
fourth order, for functions of the sixth order, I have actually expanded: but it does not
appear to contain the complete theory. Again, in the particular case of homogeneous
functions of two variables, the transforming functions may be expressed as symmetrical
functions of the roots of the equation $u = 0$, which gives rise to an entirely distinct
theory. This, however, I have not as yet developed sufficiently for publication. There
does not appear to be anything very directly analogous to the subject of this note, in
my general theory: if this be so, it proves the absolute necessity of a distinct investigation
for the present case, that which I have denominated the symmetrical one.

\bigskip
\hrule
\bigskip

% ============================================================
% PAPER 14 -- ON LINEAR TRANSFORMATIONS (start, pp.95--103)
% PDF p.117--125
% ============================================================

\begin{center}
{\Large\bfseries 14.}\\[0.6em]
{\large\bfseries ON LINEAR TRANSFORMATIONS.}\\[0.4em]
\textit{[From the Cambridge and Dublin Mathematical Journal, vol.~i.~(1846), pp.~104--122.]}
\end{center}

\medskip

In continuing my researches on the present subject, I have been led to a new
manner of considering the question, which, at the same time that it is much more
general, has the advantage of applying directly to the only case which one can
possibly hope to develope with any degree of completeness, that of functions of two
variables. In fact the question may be proposed, ``To find all the derivatives of any
number of functions, which have the property of preserving their form unaltered after
any linear transformations of the variables.'' By Derivative I understand a function
deduced in any manner whatever from the given functions, and I give the name of
Hyperdeterminant Derivative, or simply of Hyperdeterminant, to those derivatives which
have the property just enunciated. These derivatives may easily be expressed explicitly,
by means of the known method of the separation of symbols. We thus obtain the
most general expression of a hyperdeterminant. But there remains a question to be
resolved, which appears to present very great difficulties, that of determining the
independent derivatives, and the relation between these and the remaining ones. I
have only succeeded in treating a very particular case of this question, which shows
however in what way the general problem is to be attacked.

Imagine $p$ series each of $m$ variables
\[
x_1,\ y_1,\ldots\&\text{c.}\quad x_2,\ y_2,\ldots\&\text{c.}\quad \ldots\quad x_p,\ y_p,\ldots\&\text{c.},
\]
where $p$ is at least as great as $m$.

Similarly $p'$ series each of $m'$ variables
\[
x_1',\ y_1',\ldots\&\text{c.}\quad x_2',\ y_2',\ldots\&\text{c.}\quad \ldots\quad x_{p'}',\ y_{p'}',\ldots\&\text{c.},
\]
$p'$ at least as great as $m'$, and so on. Let the analogous variables $\bar{x}, \bar{y}\ldots$ be connected
with these by the equations
\[
\begin{aligned}
\bar{x} &= \lambda\,\bar{x} + \mu\,\bar{y} + \ldots,\\
\bar{y} &= \lambda'\,\bar{x} + \mu'\,\bar{y} + \ldots,\\
\bar{x}' &= \lambda\,\bar{x}' + \mu\,\bar{y}' + \ldots,\\
\bar{y}' &= \lambda''\,\bar{x}' + \mu''\,\bar{y}' + \ldots,\\
&\ \vdots
\end{aligned}
\]
where $x, y, \ldots$ stand for $x_1, y_1, \ldots$ or $x_2, y_2, \ldots$ or $x_p, y_p,\ldots$; $x', y', \ldots$ stand for $x_1', y_1',\ldots$
or $x_2', y_2',\ldots$ or $x_{p'}', y_{p'}',\ldots$ \&c. \ldots The coefficients $\lambda, \mu, \ldots, \lambda', \mu', \ldots$ \&c.; $\lambda', \mu',\ldots, \lambda'', \mu'',\ldots$
remain the same in all these systems. Suppose next,
\[
\xi = \delta_x,\quad \eta = \delta_y,\,\ldots
\]
i.e.
\[
\xi_1 = \delta_{x_1},\quad \eta_1 = \delta_{y_1},\,\ldots\,\xi_2 = \delta_{x_2},\,\ldots
\]
(where $\delta_x, \delta_y,\ldots$ are the symbols of differentiation relative to $x, y,$ \&c.). Then evidently
\[
\xi = \lambda\xi + \lambda'\eta + \ldots,
\]
\[
\eta = \mu\xi + \mu'\eta + \ldots,
\]
with similar equations for $\xi', \eta',\ldots$ Suppose
\[
\|\Omega\| = \left|\,\begin{array}{cccc} \xi_1, & \xi_2, & \ldots & \xi_p \\ \eta_1, & \eta_2, & \ldots & \eta_p \\ \multicolumn{4}{c}{\vdots} \end{array}\right|,\qquad
\|\Omega'\| = \left|\,\begin{array}{cccc} \xi_1', & \xi_2', & \ldots & \xi_{p'}' \\ \eta_1', & \eta_2', & \ldots & \eta_{p'}' \\ \multicolumn{4}{c}{\vdots} \end{array}\right|,
\]
that is to say $\|\Omega\|$ is the series of determinants formed by choosing any $m$ vertical
columns to compose a determinant, and similarly $\|\Omega'\|$, \&c. Suppose, besides,
\[
E = \left|\,\begin{array}{ccc} \lambda, & \mu, & \ldots \\ \lambda', & \mu', & \ldots \\ \multicolumn{3}{c}{\vdots} \end{array}\right|,\qquad
E' = \left|\,\begin{array}{ccc} \lambda', & \mu', & \ldots \\ \lambda'', & \mu'', & \ldots \\ \multicolumn{3}{c}{\vdots} \end{array}\right|.
\]
Then, by the known properties of determinants,
\[
\|\overline{\Omega}\| = E\,\|\Omega\|,\quad \|\overline{\Omega'}\| = E'\,\|\Omega'\|\ \&\text{c.},
\]
i.e. the terms on the one side are respectively equal to the terms on the other. Hence if
\[
\Box = F(\|\Omega\|^f,\ \|\Omega'\|^{f'},\ldots),
\]
i.e. $\Box$ a rational and integral function, homogeneous of the order $f$ in the quantities
of the series $\|\Omega\|$, homogeneous of the order $f'$ in the quantities of the series $\|\Omega'\|$,
\&c., we have immediately
\[
\overline{\Box} = E^f E'^{f'} \ldots \Box;
\]
or if $U$ be any function whatever of the variables $x, y \ldots$ which is transformed by
the linear substitutions above into $\overline{U}$, then
\[
\overline{\Box}\,\overline{U} = E^f E'^{f'} \ldots \Box\,U;
\]
or the function
\[
\Box\,U
\]
is by the above definition a hyperdeterminant derivative. The symbol $\Box$ may be called
``symbol of hyperdeterminant derivation,'' or simply ``hyperdeterminant symbol.''

Let $A, B,\ldots$ represent the different quantities of the series $\|\Omega\|$, $A', B',\ldots$ those of
the series $\|\Omega'\|$, \&c.; then $\Box$ may be reduced to a single term, and we may write
\[
\Box = A^a B^b \ldots A'^{a'} B'^{b'} \ldots
\]
Also $U$ may be supposed of the form
\[
U = \Phi \Phi' \ldots
\]
where $\Phi, \Phi'$ are functions of the variables of one of the sets $x, y,\ldots$, of one of the
sets $x'', y'',\ldots$, \&c., thus $\Phi$ is of the form
\[
F(x_1, y_1, \ldots x_1'', y_1'',\ldots),
\]
and so on. The functions $\Phi, \Phi'\ldots$ may be the same or different. It may be supposed
after the differentiations that several of the sets $x, y,\ldots$ or of the sets $x', y',\ldots$
become identical: in such cases it will always be assumed that the functions $\Phi,\ldots$
into which these sets of variables enter, are similar; so that they become absolutely
identical, when the variables they contain are made so. Thus the general expression
of a hyperdeterminant is
\[
\Box U = A^a B^b \ldots A'^{a'} B'^{b'} \ldots \Phi^k \ldots
\]
in which, after the differentiations, any number of the sets of variables are made equal.
For instance, if all the sets $x, y \ldots$ and all the sets $x', y' \ldots$ are made equal, the
hyperdeterminant refers to a single function $F(x, y \ldots x', y' \ldots)$. In any other case it
refers not to a single function but to several.

What precedes, is the general theory: it might perhaps have been made clearer
by confining it to a particular case: and by doing this from the beginning it will
be seen that it presents no real difficulties. Passing at present to some developments,
to do this, I neglect entirely the sets $x', y' \ldots$ and I assume that the number $m$ of
variables in each of the sets $x, y \ldots$ reduces itself to two; so that I consider functions
of two variables $x, y$ only. The functions $\Phi, \Phi'$, \&c. reduce themselves to functions
$V_1, V_2 \ldots V_p$ of the variables $x_1, y_1,$ or $x_2, y_2 \ldots$ or $x_p, y_p$. Writing also
\[
\Box = 12^\alpha\,13^\beta\,14^\gamma \ldots\,23^{\alpha'}\,24^{\beta'} \ldots\,34^{\alpha''} \ldots\,
\]
the symbols $A, B \ldots$ reduce themselves to 12, 13 \ldots Hence for functions of two
variables, there results the following still tolerably general form
\[
\Box U = 12^\alpha\,13^\beta\,14^\gamma \ldots\,23^{\alpha'}\,24^{\beta'} \ldots\,34^{\alpha''} \ldots\,V_1 V_2 V_3 V_4 \ldots
\]
The functions $V_1, V_2 \ldots$ may be the same or different: but they will be supposed the
same whenever the corresponding variables are made equal. This equality will be
denoted by writing, for instance,
\[
\Box V V' V V' \ldots
\]
to represent the value assumed by
\[
\Box V_1 V_2 V_3 V_4 \ldots
\]
when after the differentiations
\[
x_1 = x_3,\ y_1 = y_3,\ x_2 = x_4,\ y_2 = y_4 \,;
\]
\&c.

It is easy to determine the general term of $\Box U$. To do this, writing for shortness
\[
\alpha + \beta + \gamma \ldots = f_1,
\]
\[
\alpha + \beta' + \gamma' \ldots = f_2,
\]
\[
\beta + \beta' + \gamma'' \ldots = f_3,
\]
\[
\&\text{c.}
\]
\[
\delta_{x_1}^{f_1 - r - s - t \ldots} \delta_{y_1}^{r + s + t \ldots} V_1 = V_1{}^{r + s + t \ldots}\ \text{or}\ \delta_{x_1}^{r + s + t \ldots} \delta_{y_1}^{f_1 - r - s - t} V_1 = V_1{}^{-r - s - t \ldots}\ \text{or}\ V_1{}^{,r+s+t\ldots},
\]
the general term is
\[
\pm\,V_1^{r+s+t\ldots}\,V_2^{\alpha-r+s'+t'\ldots}\,V_3^{\beta-s+\beta'-s'+t''\ldots}\,\ldots
\]
where $r, s, t,\ldots s', t',\ldots t'',\ldots$ extend from 0 to $\alpha, \beta, \gamma,\ldots\beta', \gamma',\ldots\gamma'',\ldots$ respectively. It
would be easy to change this general term in a way similar to that which will be
employed presently for the particular case of $\Box V_1 V_2 V_3 V_4$.

If several of the functions become identical, and for these some of the letters
$f$ are equivalent, it is clear that the derivative $\Box U$ refers to a certain number of
functions $V_1, V_2 \ldots$ the same or different, of the variables $x, y; x', y';\ldots$ and besides
that this derivative is homogeneous, of the degrees $\theta_1, \theta_1',\ldots$ with respect to the
differential coefficients of the orders $f_1, f_1', \ldots$ \&c. of $V_1$, (consequently homogeneous of
the order $\theta_1 + \theta_1' + \ldots$ with respect to these differential coefficients collectively), homogeneous
and of the degrees $\theta_2, \theta_2',\ldots$ with respect to the differential coefficients of
the orders $f_2, f_2' \ldots$ of $V_2$, (consequently of the order $\theta_2 + \theta_2' \ldots$ with respect to these
collectively), and so on. The degree with respect to all the functions is of course
$\theta_1 + \theta_1' \ldots + \theta_2 + \theta_2' + \ldots,\,=p$ suppose. In general, only a single function will be considered,
and it will be assumed that $\Box U$ only contains the differential coefficients of the
$f^{\text{th}}$ order. In this case, the derivative is said to be of the degree $p$ and of the
order $f$. The most convenient classification is by degrees, rather than by orders.

Commencing with the simplest case, that of functions of the second order (and
writing $V, W$ instead of $V_1, V_2$), we have
\[
\Box V W = 12^\alpha\,V W,
\]
(where $\xi_1, \eta_1$ apply to $V$ and $\xi_2, \eta_2$ to $W$). This will be constantly represented in
the sequel by the notation
\[
12^\alpha\,V W = B_\alpha\,(V, W).
\]
Hence, writing
\[
\delta_x^a\,V = V^{,a},\quad \delta_x^{a-1} \delta_y\,V = V^{,a-1,1},\,\ldots,
\]
we have
\[
B_\alpha\,(V, W) = V^{,a}W^{a,} - {\textstyle\binom{a}{1}}\,V^{,a-1,1}W^{a-1,1} + \ldots;
\]
and in particular, according as $a$ is odd or even,
\[
B_a\,(V, V) = 0,
\]
\[
B_a\,(V, V) = V^{,a}V^{a,} - V^{a-1,1}V^{,a-1,1} + \ldots,
\]
continued to the term which contains $V^{\frac{a}{2}, \frac{a}{2}} V^{,\frac{a}{2},\frac{a}{2}}$, the coefficient of this last term
being divided by two.

Thus, for the functions $\tfrac{1}{2}(a x^2 + 2 b x y + c y^2)$, $\tfrac{1}{4}(a x^4 + 4 b x^3 y + 6 c x^2 y^2 + 4 d x y^3 + e y^4)$, \&c.,
if $a$ be made equal to 2, 4, \&c. respectively, we have the constant derivatives
\[
\begin{aligned}
ac - b^2,\\
ae - 4bd + 3c^2,\\
ag - 6bf + 15ce - 10d^2,\\
ai - 8bh + 28cg - 56df + 35e^2,
\end{aligned}
\]
which have all of them the property of remaining unaltered, \textit{à un facteur près}, when
the variables are transformed by means of $x = \lambda x' + \mu y'$, $y = \lambda' x' + \mu' y'$. Thus, for instance,
if these equations give
\[
ax^2 + 2bxy + cy^2 = \bar{a}\bar{x}^2 + 2\bar{b}\bar{x}\bar{y} + \bar{c}\bar{y}^2,
\]
then
\[
\bar{a}\bar{c} - \bar{b}^2 = (\lambda\mu' - \lambda'\mu)^2 \cdot (ac - b^2).
\]
and so on. This is the general property, which we call to mind for the case of these
constant derivatives.

The above functions may be transformed by means of the identical equation
\[
B_a\,(V, W) = 12^{a-k}\,B_k\,(V, W),
\]
to make use of which, it is only necessary to remark the general formula
\[
\delta_{x_1}^{a-k}\delta_{x_2}^{a-k}\,V W = V^{,a-k}W^{,a-k}.
\]
Thus, if $k=1$, we obtain for the above series, the new forms
\[
\begin{aligned}
ac &- b^2,\\
(ae - bd) &- 3\,(bd - c^2),\\
(ag - bf) &- 5\,(bf - ce) + 10\,(ce - d^2),\\
&\&\text{c.},\\
(ai - bh) &- 7\,(bh - cg) + 21\,(cg - df) - 35\,(df - e^2),
\end{aligned}
\]
the law of which is evident. This shows also that these functions may be linearly
expressed by means of the series of determinants
\[
\left|\,\begin{array}{cc} a, & b \\ b, & c \end{array}\right|,\quad
\left|\,\begin{array}{ccc} a, & b, & c \\ b, & c, & d \end{array}\right|,\quad \&\text{c.}
\]

We may also immediately deduce from them the derivatives $B$ which relate to two
functions. For example, for functions of the sixth order this is
\[
a g' + a' g - 6\,(b f' + b' f) + 15\,(c e' + c' e) - 20 d d',
\]
which has an obvious connection with
\[
a g - 6 b f + 15 c e - 10 d^2;
\]
and the same is the case for functions of any order.

The following theorem is easily verified; but I am unacquainted with the general
theory to which it belongs.

``If $U, V$ are any functions of the second order, and $W = \lambda U + \mu V$; then
\[
B_2'\bigl[B_2(W, W),\ B_2(W, W)\bigr] = 0
\]
(where $B_2'$ relates to $\lambda, \mu$) is the same that would be obtained by the elimination of
$x, y$ between $U=0,\ V=0$.''\footnote{Not given with the present paper.} (See Note.)

In fact this becomes
\[
4\,(ac - b^2)(a'c' - b'^2) - (ac' + a'c - 2bb')^2 = 0,
\]
which is one of the forms under which the result of the elimination of the variables
from two quadratic equations may be written. This is a result for which I am
indebted to Mr Boole.

\medskip

Passing to the third degree, we may consider in particular the derivatives
\[
\Box U V W = 23^\alpha\,31^\alpha\,12^\alpha\,U V W = C_\alpha\,(U, V, W);
\]
writing for shortness
\[
\delta_{x_2}^{\alpha-t}\,\delta_{y_2}^{t}\,\delta_{x_3}^{\alpha-r}\,\delta_{y_3}^{r}\,U = U^{-r, t},
\]
we have the general term
\[
\pm\,U^{-r, t}\,V^{-s, r}\,W^{-t, s},
\]
where $r, s, t$ extend from 0 to $\alpha$. By changing the suffixes $r, s$ the following more
convenient formula
\[
\pm\,U^{-\rho, t}\,V^{-\sigma, \rho}\,W^{-t, \sigma},
\]
where $t$ extends from 0 to $2\alpha$: $\rho, \sigma$, and $3\alpha - \rho - \sigma$ must be positive and not greater
than $2\alpha$.

In particular, according as $\alpha$ is odd or even,
\[
C_\alpha\,(U, U, U) = 0,
\]
\[
C_\alpha\,(U, U, U) = 6\Sigma\Sigma\bigl\{(-)^{t+s}\,U^{-\rho, t}\,U^{-\sigma, \rho}\,U^{-t, \sigma}\bigr\} - \tfrac{1}{2}\bigl[(-)^{\frac{a}{2}}\,U^{,\frac{a}{2},\frac{a}{2}}\bigr]^3,
\]
omitting therein those values of $\rho, \sigma$ for which $\rho > \sigma$ or $\sigma > 3\alpha - \rho - \sigma$, and dividing by
two the terms in which $\rho = \sigma$ or $\sigma = 3\alpha - \rho - \sigma$, and by six the term for which
\[
\rho = \sigma = 3\alpha - \rho - \sigma,\ = \alpha.
\]

In particular, for functions of the fourth or eighth orders we have the constant
derivatives
\[
ace - ad^2 - b^2 e - c^3 + 2 b c d;
\]
\[
aei - 8 b d - 4 a f h + 8 a g^2 + 8 i c^2 + 12 b e h - 8 c h d - 8 b g f - 22 c e g + 24 c f^2 + 24 d^2 g - 36 d e f + 15 e^3
\]
the first of which is a simple determinant. Thus we have been led to the functions
$ae - 4bd + 3c^2$ and $ace - ad^2 - eb^2 - c^3 + 2bcd$, which occur in my ``Note sur quelques
formules \&c.'' (\textit{Crelle, vol.~xxix.}~[1845] [15]), and in the forms which M. Eisenstein has
given for the solutions of equations of the first four degrees.

Let $U$ be a function of the order $4\alpha$: the derivative $C$ may be expressed by
means of the derivatives $B$.

For, consider the function
\[
B_a[U, B_a(V, W)];
\]
paying attention to the signification of $B$, this may be written
\[
1\overline{0}^{2a}\,23^{2a}\,U V W,
\]
where the symbols $\xi_0, \eta_0$ refer to the two systems $x_2, y_2$: $x_3, y_3$. Thus it is easily seen
that we may write
\[
\xi_0 = \xi_2 + \xi_3,\quad \eta_0 = \eta_2 + \eta_3,\ \text{or}\ 10 = 12 + 13 = 12 - 31,
\]
whence the function becomes
\[
(12 - 31)^{2a}\,23^{2a}\,U V W,
\]
of which all the terms vanish except
\[
12^a\,23^a\,31^a\,U V W.
\]
\[
\left[4 a\right]^2 \div \left[2 a\right]^2\,\ldots\,\left[a\right]^a
\]
Hence putting
\[
K = \frac{[2 a]^{2a}}{[2 a]^a} = \frac{2^{2a} \cdot 1 \cdot 3 \ldots (4 a - 1)}{2 \cdot 4 \ldots 4 a},
\]
we have
\[
B_a\,\bigl[U,\,B_a(V, W)\bigr] = K\,C_a\,(U, V, W),
\]
or in particular
\[
B_a\,\bigl[U,\,B_a(U, U)\bigr] = K\,C_a\,(U, U, U).
\]
Thus for example, neglecting a numerical factor,
\[
(a x^2 + 2 b x y + c y^2)\,(c x^2 + 2 d x y + e y^2) - (b x^2 + 2 c x y + d y^2)^2
\]
\[
= (ac - b^2)\,x^4 + 2\,(ad - bc)\,x^3 y + (ae + 2bd - 3c^2)\,x^2 y^2 + 2\,(be - cd)\,x y^3 + (ce - d^2)\,y^4,
\]
and then
\[
e\,(ac - b^2) - 4 d\,(ad - bc) + 6 c\,(ae + 2 b d - 3 c^2) - 4 b\,(be - cd) + a\,(ce - d^2)
\]
\[
= 3\,(ace - ad^2 - b^2 e - c^3 + 2 b c d).
\]

We have likewise the singular equation
\[
B_{2a}\,(V, W) = K\,\Bigl(\tfrac{d^a}{dx^a}\tfrac{d^a}{dy^a} - \ldots \pm \tfrac{d^{2a}}{dx^{2a}}\Bigr)\,c_a\,(U, V, W)
\]
where
\[
U = \tfrac{1}{[4a]}\,\bigl(a x^{4a} - \tfrac{[4a]}{[1]}\,b x^{4a-1} y \,\ldots\,+ z y^{4a}\bigr) \cdot \&\text{c.}
\]
If however $U = V = W$, we must write
\[
B_{2a}\,(U, U) = \tfrac{1}{6} K\,\Bigl(\tfrac{d^{2a}}{dx^{2a}}\tfrac{d^{2a}}{dy^{2a}} - \ldots \pm \tfrac{d^{4a}}{dx^{4a}}\Bigr)\,c_a\,(U, U, U),
\]
the reason of which is easily seen. This subject will be resumed in the sequel.

The functions $C$ may be transformed in the same way as the functions $B$ have been.
In fact
\[
C_a\,(U, V, W) = 12^{a-k}\,23^{a-k}\,31^{a-k}\,C_k\,(U, V, W);
\]
if in particular $k=1$, then
\[
\ldots U^{,a-k}\ \text{for}\ U^{,a-k},\,\&\text{c}.
\]
but in general
\[
\xi_1^{,\rho}\,\eta_1^{,\sigma}\,\xi_2^{,\rho'}\,\eta_2^{,\sigma'}\,\xi_3^{,t}\,\eta_3^{,t'}\,C_k\,(U, V, W),\ \text{where}\ \rho + \rho' = \sigma + \sigma' = t + t' = 2a - 2,
\]
\[
2a - 2,\quad 4a - 4,\quad 2a - 2,\quad U^{,\rho}\,V^{,\sigma}\,W^{,t}\ \text{for}\ U^{,\rho},\ \&\text{c.}
\]
whence
\[
C_a\,(U, V, W) = \Sigma\Sigma\bigl\{(-)^{r+t}\,U^{,\rho}\,V^{,\rho+1}\,V^{,\sigma}\,W^{,3a-\rho-\sigma-1}\bigr\}\cdot\Sigma\bigl[(-)^t\,A_r{}^t\,A_{\sigma+1-r}\,A_t\bigr],
\]
where $A_p{}^t = \tfrac{[a-1]^p}{[t]^p}$; $t$ extends from 0 to $a-1$; $\rho, \sigma - 1$, and $3a - \rho - \sigma - 2$ may have
each of them any positive values not greater than $2a - 2$.

In particular
\[
C_a\,(U, U, U) = 6\,\Sigma\Sigma\bigl\{(-)^{r+t}\,U^{,\rho}\,U^{,\sigma-1}\,U^{,3a-\rho-\sigma-2}\bigr\}\cdot\Sigma\bigl[(-)^{r+t}\,A_r{}^t\,A_{\sigma+1-r}\,A_t\bigr],
\]
where $\rho, \sigma$ need only have such values that $\rho < \sigma - 1, \sigma - 1 < 3a - \rho - \sigma - 2$.

In particular the derivative $aei - \ldots + 15 e^3$ may be transformed into
\[
\left|\,\begin{array}{ccc} a, & d, & g \\ b, & e, & h \\ c, & f, & i \end{array}\right|
- \left|\,\begin{array}{ccc} a, & e, & f \\ b, & f, & g \\ c, & g, & h \end{array}\right|
- \left|\,\begin{array}{ccc} b, & c, & g \\ c, & d, & h \\ d, & e, & i \end{array}\right|
+ 6 \left|\,\begin{array}{ccc} b, & d, & f \\ c, & e, & g \\ d, & f, & h \end{array}\right|
- 15 \left|\,\begin{array}{ccc} c, & d, & e \\ d, & e, & f \\ e, & f, & g \end{array}\right|
\]
in which form it is obviously a linear function of the determinants

\end{document}
